% Generated from locale/tr/content/incompleteness/introduction/historical-background.tex; do not hand-edit.


\olfileid[tr]{inc}{int}{bgr}

\olsection{Tarihsel Arka Plan}

Bu bölümde eksiklik teoremlerini bağlamına yerleştirmeye yardımcı olacak
tarihsel gelişmeleri kısaca ele alacağız. Özellikle matematiksel mantık
tarihinin çok kaba bir genel görünümünü verecek, ardından matematiğin
temellerinin tarihi hakkında birkaç söz söyleyeceğiz.

\begin{digress}
  ``Matematiksel mantık'' sözü iki anlamlıdır. ``Matematiksel'' sözcüğü,
  ``matematiğin mantığı''nda olduğu gibi konuyu niteliyor ve matematiksel
  akıl yürütmenin ilkelerini belirtiyor diye yorumlanabilir; yahut
  ``mantığın matematiği''nde olduğu gibi yöntemleri niteliyor ve akıl
  yürütme ilkelerinin matematiksel incelenmesini belirtiyor diye
  yorumlanabilir. Aşağıdaki anlatı, çoğu kez aynı anda olmak üzere,
  matematiksel mantığı her iki anlamıyla da içerir.
\end{digress}

Mantık incelemesi esas olarak yaklaşık 384--322 \textsc{MÖ} yıllarında yaşamış olan
Aristoteles'le başladı. Onun \emph{Kategoriler}, \emph{Birinci Analitikler}
ve \emph{İkinci Analitikler} adlı eserleri, kıyasın kapsamlı ve sistemli bir
incelemesi de dâhil olmak üzere, bilimsel akıl yürütme ilkelerinin sistemli
incelemelerini içerir.

Aristoteles'in mantığı Orta Çağ boyunca skolastik felsefeye egemen oldu;
hatta on sekizinci yüzyılın sonlarında Kant, Aristoteles mantığının kusursuz
olduğunu ve gözden geçirilmeye ihtiyaç duymadığını savunuyordu. Ne var ki
kıyas kuramı, matematiksel akıl yürütmenin en yüzeysel yönlerinden başka bir
şeyi modellemek için fazlasıyla sınırlıdır. Bir yüzyıl önce Newton'un çağdaşı
Leibniz, mantıksal akıl yürütme için eksiksiz bir ``hesap'' tasarlamış ve
özünde önerme mantığının bir sürümünü betimleyerek böyle bir hesabı kurma
yolunda bazı ilk adımları atmıştı.

On dokuzuncu yüzyıl mantık için bir dönüm noktasıydı. George Boole 1854'te,
çağdaş sunumlardan çok uzak olmayan kapsamlı bir cebirsel önerme mantığı
incelemesi olan \emph{Düşüncenin Yasaları}'nı yazdı. Gottlob Frege 1879'da,
önerme mantığını niceleyiciler ve bağıntılarla genişleten ve böylece birinci
derece mantığı içeren \emph{Begriffsschrift}'ini (Kavram Yazısı) yayımladı.
Gerçekte Frege'nin mantıksal sistemleri daha yüksek dereceli mantığı ve
başka şeyleri de içeriyordu. Frege \emph{Aritmetiğin Temel Yasaları}'nda,
bütün aritmetiğin kendi Begriffsschrift'inde salt mantıksal varsayımlardan
türetilebileceğini göstermeye girişti. Ne yazık ki Russell'ın 1902'de
gösterdiği üzere bu varsayımlar tutarsız çıktı. Ancak tutarsız aksiyom bir
yana bırakılırsa Frege çağdaş mantığı aşağı yukarı tek başına icat etmişti;
bu şaşırtıcı bir başarıdır. Niceleme mantığı, Boole'dan sonra Peirce ve
Schr\"oder'in de aralarında bulunduğu cebir eğilimli düşünürlerce bağımsız
olarak geliştirildi.

Şimdi matematiğin temellerindeki gelişmelere dönelim. Mantık matematikte
önemli bir rol oynadığından, elbette az önce betimlenen gelişmelerle büyük
bir etkileşim vardır. Örneğin Frege mantığını, bütün matematiğin yalnızca
mantıksal çerçevesine dayandırılabileceğini gösterme açık amacıyla
geliştirdi; özellikle matematiğin, Kant'ın ileri sürdüğü gibi a priori
\emph{sentetik} doğrulardan değil a priori \emph{analitik} doğrulardan
oluştuğunu göstermek istiyordu.

Birçok kişi asıl matematiğin doğuşunu Yunanlarla başlatır. Yaklaşık MÖ
300'de yazılan Öklid'in \emph{Ögeler}'i, kesinlik ve titizliğe yaptığı
vurgu ile Yunan matematiğinin olgun bir örneğidir. Öklid'in \emph{Ögeler}'indeki
tanımlar ve ispatlar bugün lise geometri kitaplarında aşağı yukarı olduğu
gibi yaşamaktadır (geometrinin liselerde hâlâ öğretildiği ölçüde). Bu
matematiksel akıl yürütme modeli yalnızca matematikte değil, felsefenin
dallarında da titiz kanıtlamanın paradigması sayılmıştır. (Spinoza ahlaki ve
dinsel savları bile Öklid tarzında sunmuştur; bunu görmek tuhaftır!)

Kalkülüs on yedinci yüzyılda Newton ve Leibniz tarafından icat edildi.
(Öncelik konusunda yüzyıllarca şiddetli bir tartışma sürdü; fakat bugün çoğu
araştırmacı iki gelişmenin büyük ölçüde bağımsız olduğu kanısındadır.)
Kalkülüs, örneğin sonsuz küçük niceliklerin sonsuz toplamları üzerine akıl
yürütmeyi içerir; bu özellikler, Tanrı'ya inanmanın kendi çağının
matematiğinden daha az akılcı olmadığını savunan Piskopos Berkeley'nin
eleştirilerini besledi. Kalkülüs yöntemleri on sekizinci yüzyılda yaygın
biçimde kullanıldı; örneğin Leonhard Euler sonsuz toplamları içeren
hesaplarla çarpıcı sonuçlar elde etti.

On dokuzuncu yüzyılda matematikçiler kalkülüsü daha sağlam bir temele
oturtarak Berkeley'nin eleştirilerini karşılamaya çalıştılar. Cauchy,
Weierstrass, Bolzano ve başkalarının çabaları limit, süreklilik, türev ve
integral için günümüzdeki ``epsilon ve delta'' tanımlarına, başka bir
deyişle sonsuz küçüklere hiç başvurmayan tanımlara yol açtı. Yüzyılın daha
sonraki döneminde matematikçiler daha ileri gitmeye ve reel sayıların
kendileri de dâhil olmak üzere kalkülüsün bütün yönlerini doğal sayılar
cinsinden açıklamaya çalıştılar. (Kronecker: ``Tam sayıları Tanrı yarattı;
geri kalan her şey insanın eseridir.'') Dedekind 1872'de ``Süreklilik ve
irrasyonel sayılar''ı yazdı ve reel sayıların rasyonel sayı kümeleri olarak
(bunlar da bildiğiniz gibi doğal sayı çiftleri olarak görülebilir) nasıl
``kurulacağını'' gösterdi; 1888'de doğal sayıları salt ``mantıksal''
terimlerle açıklamayı amaçlayan ``Was sind und was sollen die Zahlen''ı
(kabaca, ``Doğal sayılar nedir ve ne olmalıdır?'') yazdı. Kronecker 1887'de
bütün matematiksel nesnelerin tam sayılar cinsinden temsilinden söz ettiği
``\"Uber den Zahlbegriff''i (``Sayı kavramı üzerine'') yazdı; Giuseppe Peano
1889'da doğal sayılar için biçimsel, simgesel aksiyomlar verdi.

On dokuzuncu yüzyılın sonu sonsuzla uğraşmada yeni bir cesaret de getirdi.
O zamana kadar sonsuz nesneler ve yapılar (doğal sayılar kümesi gibi)
çekingenlikle ele alınıyor; ``sonsuz sayıda'', ``istediğiniz kadar çok'' ve
``limitte yaklaşır'', ``istediğiniz kadar yakınlaşır'' diye anlaşılıyordu.
Georg Cantor ise sonsuzu olduğu gibi ele almanın mümkün olduğunu gösterdi.
Cantor, Dedekind ve başkalarının çalışmaları, bugün yaygın biçimde kabul
edilen genel küme-kuramsal matematik anlayışının ortaya çıkmasına yardım
etti.

Bu bizi yirminci yüzyılda mantık ve temellerdeki gelişmelere getirir.
Russell 1902'de Frege'nin mantıksal sistemindeki paradoksu keşfetti. Zermelo
1904'te ``seçim aksiyomu'' denen ilkeyi kullanarak Cantor'un iyi sıralama
ilkesini ispatladı; bu aksiyomun meşruluğu büyük tartışma yarattı. Russell
ile Whitehead'in, matematiği mantıksal temeller üzerinde kurmaya yönelik
Fregeci programı genişleten üç ciltlik \emph{Principia Mathematica}'sı
1910 ile 1913 arasında yayımlandı. Ne yazık ki Russell ile Whitehead'in salt
mantıksal olarak gerekçelendirilmesi zor görünen iki ilkeyi benimsemesi
gerekti: sonsuzluk aksiyomu ve ``indirgenebilirlik'' aksiyomu. Poincar\'e
1900'lü yıllarda matematikte ``yüklemsel olmayan tanımlar''ın kullanılmasını
eleştirdi; Brouwer ise 1910'lu yıllarda dışlanmış orta yasasını ($!A \lor
\lnot !A$) kullanmaktan kaçınan ``sezgici'' bir temel üzerinde bütün
matematiği yeniden kurmayı önermeye başladı.

Gerçekten tuhaf günlerdi!{} Bütün matematiği mantığa indirgeme programına
bugün ``mantıkçılık'' denir ve yukarıda anılan güçlükler yüzünden genellikle
başarısız olduğu düşünülür. Matematiği sezgici zihinsel kuruluşlar cinsinden
geliştirme programına ``sezgicilik'' denir ve gündelik matematiğe aşırı ağır
kısıtlamalar getirdiği kabul edilir. Yüzyılın başlarında tüm zamanların en
etkili matematikçilerinden David Hilbert, Cantor ve Dedekind'in getirdiği
yeni, soyut yöntemlerin güçlü bir savunucusuydu: ``Cantor'un bizim için
yarattığı cennetten kimse bizi kovamayacaktır.'' Aynı zamanda bu yeni
(gariptir ki artık ``klasik'' denen) yöntemlere yöneltilen temel eleştirilerine
duyarlıydı. Hem pastaya sahip olup hem onu yemeyi sağlayacak bir yol önerdi:
\begin{enumerate}
\item Klasik yöntemleri biçimsel aksiyomlar ve kurallarla temsil etmek;
  matematiksel soruları aksiyomatik bir sistemde !!{formula}s olarak
  temsil etmek.
\item Bu biçimsel tümdengelim sistemlerinin tutarlı olduğunu güvenli,
  ``sonlu'' yöntemlerle ispatlamak.
\end{enumerate}

Hilbert'in çalışması ilk hedefi gerçekleştirme yolunda büyük mesafe aldı.
Ünlü \emph{Geometrinin Temelleri} kitabında bunu 1899'da geometri için
yapmıştı. Sonraki yıllarda kendisi, öğrencilerinin ve çalışma arkadaşlarının
bir bölümüyle birlikte, matematiğin başka alanlarında Hilbert'in geometri
için yaptığını yapmaya çalıştı. Hilbert aritmetik ve analiz için aksiyom
sistemleri verdi. Zermelo küme kuramını aksiyomlaştırdı; bu aksiyomlaştırma
Fraenkel, Skolem, von Neumann ve başkalarınca genişletildi. 1920'lerin
ortasına gelindiğinde ``bütün'' matematiğin aksiyomlaştırılması unvanında hak
iddia eden iki yaklaşım vardı: Russell ile Whitehead'in \emph{Principia
mathematica}'sı ve Zermelo--Fraenkel küme kuramı diye tanınan yaklaşım.

Hilbert 1921'de bu sistemlerin tutarlı olduğunu ispatlama hedefiyle bir
araştırma programı başlattı. Bu projede birkaç öğrencisi, özellikle Bernays,
Ackermann ve daha sonra Gentzen kendisine yardım etti. Hedefe ulaşmanın
temel düşüncesi, matematikteki bir tutarsızlığın !!{derivation} olanağı
sorusunu muhtemel simge dizileri hakkında birleşimsel bir problem olarak
kurmaktı: sözgelimi aritmetik, analiz ya da küme kuramının aksiyomatik bir
sisteminin aksiyomlarından $!A \land \lnot !A$'nın doğru bir
!!{derivation}{}i olma ölçütünü karşılayan muhtemel tümce dizileri. Böyle bir
simge dizisinin olanaksızlığının ispatı, kendisi de matematiksel bir ispat
olduğu için, bu aksiyomatik sistemlerde biçimselleştirilebilir olacaktı.
Başka bir deyişle, örneğin aritmetiğin tutarlı olduğunu söyleyen bir
$\OCon$ tümcesi olacaktı. Üstelik bu tümce, özellikle ispatı yalnızca çok
sınırlı ``sonlu'' araçlar gerektiriyorsa, söz konusu sistemlerde
ispatlanabilir olmalıydı.

Geliştirilen aksiyom sistemlerinin her matematiksel soruyu çözmesi olan
ikinci amaç iki biçimde kesinleştirilebilir. Birincisi şöyle dile
getirilebilir: Matematik için bir aksiyom sisteminin dilindeki her $!A$
tümcesi için ya $!A$ ya da $\lnot !A$ aksiyomlardan ispatlanabilir. Bu doğru
olsaydı aksiyomlar temelinde ne ispatlanabilen ne de çürütülebilen tümceler,
aksiyomların karara bağlamadığı sorular bulunmazdı. Bu özelliği taşıyan
aksiyom sistemine \emph{tam} denir. Elbette belirli bir tümce için iki
seçenekten hangisinin geçerli olduğunu belirlemek yine de güç olabilir.
Fakat ilkece bunu yapacak bir yöntem bulunmalıdır. Gerçekte Hilbert'in ele
aldığı aksiyom ve !!{derivation} sistemlerinde tamlık böyle bir yöntemin
varlığını gerektirirdi; gerçi Hilbert bunun farkında değildi. Soruyu
yorumlamanın ikinci yolu daha güçlü şu koşul olurdu: verilen bir $!A$
tümcesinin aksiyomlardan türetilebilir olup olmadığını belirleyen mekanik,
hesaplamalı bir yöntemin bulunması.

G\"odel 1931'de bu programın başarıya ulaşamayacağını gösteren iki
``eksiklik teoremi''ni ispatladı. Matematik için tam bir aksiyom sistemi
yoktur; özellikle aksiyomların tutarlılığını dile getiren tümce ne
ispatlanabilir ne de çürütülebilir bir tümcedir.

Bu, Hilbert'in özgün programına öldürücü bir darbe indirdi. Bununla birlikte,
matematikte sık sık olduğu gibi, heyecan verici yeni araştırma yolları da
açtı. Matematiğin her şeyi kapsayan tek bir biçimsel sistemi yoksa daha dar
kapsamlı sistemler geliştirip bunlarda nelerin ispatlanabileceğini
araştırmak anlamlıdır. Bu sistemlerin tutarlılığını kurmak için daha az
kısıtlı ispat yöntemleri geliştirmek ve tutarlılıklarını ispatlamanın ne
kadar güç olduğunu ölçme yolları bulmak da anlamlıdır. G\"odel (hemen
hemen) her biçimsel sistemin karara bağlayamadığı soruları olduğunu
gösterdiğine göre, belirli bir biçimsel sistemin karara bağlayamadığı
``ilginç'' soruları aramak ve bunları karara bağlamak için biçimsel bir
sistemin ne kadar güçlü olması gerektiğini belirlemek anlamlıdır.
Mantıkçılar günümüze dek bu soruları yeni bir matematiksel disiplin olan
ispat kuramı içinde araştırmayı sürdürmüşlerdir.

